\documentclass[10pt,a4paper]{article}
\usepackage[hmargin=18mm,vmargin=14mm]{geometry}
\usepackage{amsmath,amssymb}
\usepackage{booktabs,longtable,array}
\usepackage[table]{xcolor}
\usepackage{enumitem}
\usepackage[hidelinks]{hyperref}
\setlength{\parindent}{0pt}
\setlength{\parskip}{0.5em}
\newcommand{\done}[1]{\textcolor{black!45}{#1}}
\newcommand{\hw}[1]{\textbf{#1}}

\begin{document}

{\LARGE\bfseries Advanced Optimization --- Lecture Plan}\\[0.2em]
{\large Fall 2026, Sogang University, Dept.\ of Electronic Engineering \quad Prof.\ Hongseok Kim}\\[0.2em]
{\small Text: Nocedal \& Wright, \emph{Numerical Optimization}, 2nd ed.\quad Tue/Thu.\quad Revised Oct 9, 2026, v3
(student presentations Nov 10--17; Ch.~10 if time permits).}

\textbf{Scope.} The course covers \emph{unconstrained} optimization in depth (Ch.~3--7) and the \emph{theory} of
constrained optimization (Ch.~12). It closes with a two-lecture reading guide to primal-dual interior-point
methods (Ch.~19) --- the algorithm inside IPOPT and KNITRO --- and a final project in which you run IPOPT on a
small AC OPF and explain every line of its log using Ch.~3, 12 and 19. Constrained \emph{algorithms}
(Ch.~15--18) are outside the course; the last section says when you would need them.

\renewcommand{\arraystretch}{1.15}
\begin{longtable}{@{}>{\raggedright\arraybackslash}p{7mm}>{\raggedright\arraybackslash}p{18mm}>{\raggedright\arraybackslash}p{86mm}>{\raggedright\arraybackslash}p{47mm}@{}}
\toprule
\textbf{Wk} & \textbf{Date} & \textbf{Topic} & \textbf{Reading / HW}\\
\midrule
\endhead
\done{1--4} & \done{Sep 1--22} & \done{Ch.~1--2; Ch.~3 line search (Wolfe, Zoutendijk, Newton/quasi-Newton rates,
  Hessian modification); Ch.~4 framework, Algorithm 4.1, Theorem 4.1 --- \emph{done}} & \done{HW3 out Sep 22}\\
4 & Sep 24 & \emph{No class (Chuseok)} & \\
\midrule
\done{5} & \done{Sep 29} & \done{Ch.~4: Cauchy point, fraction of Cauchy decrease, dogleg, 2D subspace
  minimization; Steihaug-CG named only --- \emph{done}} & \done{\S4.1}\\
 & \done{Oct 1} & \done{Ch.~4: global convergence (Thm.~4.5--4.6), local convergence, scaling; \S4.3 skipped
  --- \emph{done}} & \done{\S4.2, 4.4--4.5; HW4 out}\\
\done{6} & \done{Oct 6} & \done{Ch.~5: linear CG I --- conjugate directions, expanding subspace (Thm.~5.2),
  Algorithm 5.1 $\to$ 5.2 --- \emph{done}} & \done{\S5.1}\\
 & \done{Oct 8} & \done{Ch.~5: linear CG II --- rate of convergence, eigenvalue clustering, preconditioning;
  nonlinear CG --- \emph{done}} & \done{\S5.1--5.2; HW5 out}\\
7 & Oct 13 & Ch.~6: BFGS --- secant equation, why $s^{T}y>0$ needs Wolfe, implementation & \S6.1\\
 & Oct 15 & Ch.~6: SR1, Broyden class, convergence results (stated), Dennis--Mor\'e revisited
  & \S6.2--6.4; \hw{HW6 out}\\
\midrule
8 & Oct 20 & \textbf{Midterm exam} (Ch.~3--6) & \\
 & Oct 22 & Buffer / catch-up (midterm week). If Ch.~6 is on schedule: \textbf{Ch.~10} Gauss--Newton and Levenberg--Marquardt & \S10.1--10.3 (if time permits)\\
\midrule
9 & Oct 27 & Ch.~7: inexact Newton, Newton-CG, trust-region Newton-CG & \S7.1\\
 & Oct 29 & Ch.~7: L-BFGS; what is skipped in \S7.3--7.4 and why & \S7.2; \hw{HW7 out}\\
10 & Nov 3 & Ch.~12: examples, tangent cone, LICQ, first-order necessary conditions (statement)
  & \S12.1--12.2\\
 & Nov 5 & Ch.~12: proof of KKT via Farkas' lemma; constraint qualifications & \S12.3, 12.6\\
\midrule
11 & Nov 10 & \textbf{Student presentation 1: Shampoo} (instructor traveling; recorded) & see below\\
 & Nov 12 & \textbf{Student presentation 2: SOAP} (instructor traveling; recorded) & see below\\
12 & Nov 17 & \textbf{Student presentation 3: Muon} (instructor traveling; recorded) & see below\\
 & Nov 19 & Ch.~12: second-order conditions, critical cone & \S12.5\\
13 & Nov 24 & Ch.~12: degeneracy, duality, multipliers as sensitivities; KKT of a small AC OPF,
  multipliers as prices & \S12.7--12.9, notes; \hw{HW12 out}\\
 & Nov 26 & Ch.~19 reading guide I: barrier problem, primal-dual Newton system, central path,
  inertia (\S3.4 again) & \S19.1--19.3\\
14 & Dec 1 & Ch.~19 reading guide II: merit function / filter line search; an IPOPT log, line by line
  & \S19.4--19.6; W\"achter--Biegler (2006)\\
 & Dec 3 & Final project workshop (IPOPT via JuMP / Pyomo / CasADi) & \\
15 & Dec 8, 10 & Final project presentations & \\
16 & Dec 15, 17 & Final exam week --- no lecture; \textbf{project report due} & \\
\bottomrule
\end{longtable}

{\small Assumed: 16-week semester from Sep 1, midterm week 8 (Oct 20--22), final week 16 (Dec 15--17).
Adjust if the official calendar differs. Changes from the Sep 22 version: the instructor is traveling
Nov 9--18, so Nov 10, 12 and 17 are student presentations; Ch.~12 continues on Nov 19--24, and Ch.~19 and the
project workshop move later by one session. The two buffer sessions (Nov 19, Dec 3) are used for this.}

\textbf{Student presentations (Nov 10, 12, 17).} Six students, three groups of two, one optimizer per session:
\begin{itemize}[leftmargin=*,itemsep=0.1em]
  \item \textbf{Shampoo} --- Gupta, Koren \& Singer (2018).
  \item \textbf{SOAP} --- Vyas et al.\ (2024).
  \item \textbf{Muon} --- Jordan et al.\ (2024).
\end{itemize}
Each group studies its method together and presents it in the language of this course: which matrix plays the
role of the preconditioner or of $B_k$, how it is updated, what one step costs, and how it relates to
preconditioning (Ch.~5) and quasi-Newton updating (Ch.~6--7). The session is run by the students and
recorded; the group uploads the recording.

\textbf{Homework thread.} One problem set per chapter, one week each. Code is expected to be written with AI
assistance; grading is on the \emph{observations}, not the code. The line search from HW3 is reused in HW5--HW7.
\begin{itemize}[leftmargin=*,itemsep=0.1em]
  \item \textbf{HW3} --- backtracking (Alg.~3.1) and strong-Wolfe (Alg.~3.5--3.6) line searches; steepest descent
    and Newton on Rosenbrock; Alg.~3.3 on an indefinite Hessian; Zoutendijk annotated; Ex.~3.5.
  \item \textbf{HW4} --- Algorithm 4.1 with dogleg on Rosenbrock; log $\Delta_k,\rho_k$; compare iteration counts
    with HW3.
  \item \textbf{HW5} --- linear CG on an SPD matrix with clustered vs.\ spread eigenvalues ($n=100$): iteration
    count follows the clusters; FR and PR$^{+}$ with the HW3 line search, $c_2=0.1$.
  \item \textbf{HW6} --- BFGS with the HW3 line search on Rosenbrock; log $\|B_k-\nabla^2 f_k\|$ against
    $\|(B_k-\nabla^2 f_k)p_k\|/\|p_k\|$ (Dennis--Mor\'e in action).
  \item \textbf{HW7} --- L-BFGS with $m\in\{3,5,20\}$ on an $n=10^4$ problem; compare with Newton-CG.
  \item \textbf{HW12} --- KKT conditions by hand for a 2- or 3-bus DC OPF; identify what each multiplier means.
\end{itemize}

\textbf{Final project.} Solve a small AC OPF (pglib-opf \texttt{case5} or \texttt{case14}) with IPOPT. Annotate
the solver log: barrier parameter $\mu$ and its decrease (\S19.3), primal/dual infeasibility and complementarity
(Ch.~12), \texttt{inertia correction} (\S3.4), step lengths $\alpha$ and the merit/filter decision (\S19.4).
Report: 4 pages, plus a 10-minute presentation.

\textbf{If time permits.} Ch.~10 (nonlinear least squares: Gauss--Newton, Levenberg--Marquardt). It needs
only Ch.~4 and Ch.~6: Gauss--Newton chooses $B_k=J_k^{T}J_k$, and Levenberg--Marquardt is the trust-region
method of Ch.~4 applied to this model. It is taught in a free session before Ch.~12 (first choice: Oct 22).

\textbf{Not covered, and when you would need it.}
\begin{itemize}[leftmargin=*,itemsep=0.1em]
  \item Ch.~8 (derivatives, automatic differentiation) --- this is what \texttt{autograd} does; read \S8.2 once.
  \item Ch.~9 (derivative-free), Ch.~11 (nonlinear equations) --- niche for this course.
  \item Ch.~13 (simplex), Ch.~14 (LP interior point) --- Ch.~14's idea is all that Ch.~19 needs, and it is given
    in the reading guide.
  \item Ch.~15--18 (merit functions, QP, penalty/augmented Lagrangian, SQP) --- read Ch.~17 if you use
    ADMM-type methods; Ch.~18 is KNITRO's other mode.
\end{itemize}

\end{document}
